
\documentclass{article}
\usepackage{colortbl}
\usepackage{makecell}
\usepackage{multirow}
\usepackage{supertabular}

\begin{document}

\newcounter{utterance}

\centering \large Interaction Transcript for game `ifeval', experiment `ifeval', episode 449 with selfrefine{-}llama{-}3.1{-}8b{-}instruct.
\vspace{24pt}

{ \footnotesize  \setcounter{utterance}{1}
\setlength{\tabcolsep}{0pt}
\begin{supertabular}{c@{$\;$}|p{.15\linewidth}@{}p{.15\linewidth}p{.15\linewidth}p{.15\linewidth}p{.15\linewidth}p{.15\linewidth}}
    \# & \multicolumn{2}{c}{Player} && \multicolumn{2}{c}{Game Master} \\
    \hline

    \theutterance \stepcounter{utterance}  
    & & & \multicolumn{4}{p{0.6\linewidth}}{
        \cellcolor[rgb]{0.9,0.9,0.9}{
            \makecell[{{p{\linewidth}}}]{
                \texttt{\tiny{[P1$\langle$GM]}}
                \texttt{I need to write a research proposal for a project on the effects of social media on the mental health of teenagers. Please only write short sentences, and don't use any commas in your entire response. Include at least five placeholder represented by square brackets, such as {[}website{]}.} \\
            }
        }
    }
     \\ \\

    \theutterance \stepcounter{utterance}  
    & & \multicolumn{4}{p{0.6\linewidth}}{
        \cellcolor[rgb]{0.95,0.95,0.95}{
            \makecell[{{p{\linewidth}}}]{
                \texttt{\tiny{[P1$\leftrightarrow$P]}}
                \texttt{{[}DRAFT{]}} \\
\texttt{Title: Exploring the Impact of Social Media on Teenagers' Mental Health} \\
\\ 
\texttt{Background: Social media has become an integral part of modern life with {[}number{]} of teenagers using platforms like {[}Facebook{]} and {[}Instagram{]} daily.} \\
\\ 
\texttt{Objectives:} \\
\texttt{{-} Investigate the relationship between social media usage and symptoms of depression and anxiety in teenagers.} \\
\texttt{{-} Examine the role of social comparison and cyberbullying in contributing to mental health issues.} \\
\texttt{{-} Identify the most commonly used social media platforms among teenagers and their associated risks.} \\
\\ 
\texttt{Methodology:} \\
\texttt{{-} Conduct a mixed{-}methods study combining both qualitative and quantitative data collection methods.} \\
\texttt{{-} Recruit {[}number{]} of participants aged {[}age range{]} from {[}school/college{]} and {[}community{]} settings.} \\
\texttt{{-} Use {[}software{]} to analyze data from social media usage and mental health surveys.} \\
\\ 
\texttt{Expected Outcomes:} \\
\texttt{{-} Quantitative data will provide insights into the prevalence of mental health issues among teenagers and their correlation with social media usage.} \\
\texttt{{-} Qualitative data will offer a deeper understanding of the experiences and perceptions of teenagers regarding social media and mental health.} \\
\\ 
\texttt{Timeline:} \\
\texttt{{-} Literature review: {[}timeframe{]}} \\
\texttt{{-} Data collection: {[}timeframe{]}} \\
\texttt{{-} Data analysis: {[}timeframe{]}} \\
\texttt{{-} Writing and submission of the thesis: {[}timeframe{]}} \\
\\ 
\texttt{Resources:} \\
\texttt{{-} {[}Funding agency{]} will provide {[}amount{]} of funding for the project.} \\
\texttt{{-} {[}University{]} will offer {[}resources{]} including {[}equipment{]} and {[}expertise{]}.} \\
\\ 
\texttt{Expected Impact:} \\
\texttt{{-} The study will contribute to the existing body of research on the effects of social media on mental health.} \\
\texttt{{-} The findings will inform the development of effective interventions and policies to mitigate the negative effects of social media on teenagers' mental health.} \\
\\ 
\texttt{Limitations:} \\
\texttt{{-} The study may be limited by the sample size and the generalizability of the findings to other populations.} \\
\\ 
\texttt{Ethics:} \\
\texttt{{-} The study will be conducted in accordance with {[}institutional review board{]} guidelines and regulations.} \\
\texttt{{-} Participants will be informed about the study's purpose and procedures and will provide consent before participating.} \\
            }
        }
    }
    & \\ \\

    \theutterance \stepcounter{utterance}  
    & & \multicolumn{4}{p{0.6\linewidth}}{
        \cellcolor[rgb]{0.95,0.95,0.95}{
            \makecell[{{p{\linewidth}}}]{
                \texttt{\tiny{[P1$\leftrightarrow$P]}}
                \texttt{{[}SELF{-}FEEDBACK{]}} \\
\texttt{FORMAT: ok} \\
\texttt{RULES: ok} \\
\texttt{STRATEGY: problem} \\
\texttt{REASON: The answer is too long and does not follow the constraint of short sentences without commas.} \\
            }
        }
    }
    & \\ \\

    \theutterance \stepcounter{utterance}  
    & \multicolumn{4}{p{0.6\linewidth}}{
        \cellcolor[rgb]{0.561,0.737,0.561}{
            \makecell[{{p{\linewidth}}}]{
                \texttt{\tiny{[P1$\rangle$GM]}}
                \texttt{Title: Exploring the Impact of Social Media on Teenagers' Mental Health} \\
\\ 
\texttt{Background: Social media is used by {[}number{]} of teenagers daily on platforms like {[}Facebook{]} and {[}Instagram{]}.} \\
\\ 
\texttt{Objectives:} \\
\texttt{{-} Investigate social media usage and symptoms of depression and anxiety.} \\
\texttt{{-} Examine social comparison and cyberbullying's role in mental health issues.} \\
\texttt{{-} Identify most used social media platforms and associated risks.} \\
\\ 
\texttt{Methodology:} \\
\texttt{{-} Conduct mixed{-}methods study combining qualitative and quantitative data.} \\
\texttt{{-} Recruit {[}number{]} participants aged {[}age range{]} from {[}school/college{]} and {[}community{]}.} \\
\texttt{{-} Use {[}software{]} to analyze data from social media usage and mental health surveys.} \\
\\ 
\texttt{Expected Outcomes:} \\
\texttt{{-} Quantitative data will show prevalence of mental health issues and correlation with social media usage.} \\
\texttt{{-} Qualitative data will offer deeper understanding of teenagers' experiences and perceptions.} \\
\\ 
\texttt{Timeline:} \\
\texttt{{-} Literature review will take {[}timeframe{]} months.} \\
\texttt{{-} Data collection will take {[}timeframe{]} months.} \\
\texttt{{-} Data analysis will take {[}timeframe{]} months.} \\
\texttt{{-} Writing and submission will take {[}timeframe{]} months.} \\
\\ 
\texttt{Resources:} \\
\texttt{{-} {[}Funding agency{]} will provide {[}amount{]} funding.} \\
\texttt{{-} {[}University{]} will offer {[}resources{]} including {[}equipment{]} and {[}expertise{]}.} \\
\\ 
\texttt{Expected Impact:} \\
\texttt{{-} Study will contribute to existing research on social media's effects on mental health.} \\
\texttt{{-} Findings will inform development of effective interventions and policies.} \\
\\ 
\texttt{Limitations:} \\
\texttt{{-} Study may be limited by sample size and generalizability of findings.} \\
\\ 
\texttt{Ethics:} \\
\texttt{{-} Study will follow {[}institutional review board{]} guidelines.} \\
\texttt{{-} Participants will be informed and provide consent before participating.} \\
            }
        }
    }
    & & \\ \\

    \theutterance \stepcounter{utterance}  
    & & & \multicolumn{2}{p{0.3\linewidth}}{
        \cellcolor[rgb]{0.9,0.9,0.9}{
            \makecell[{{p{\linewidth}}}]{
                \texttt{\tiny{[GM$|$GM]}}
                {'punctuation:no_comma': {}, 'detectable_content:number_placeholders': {'num_placeholders': 5}}
            }
        }
    }
    & & \\ \\

    \theutterance \stepcounter{utterance}  
    & & & \multicolumn{2}{p{0.3\linewidth}}{
        \cellcolor[rgb]{0.9,0.9,0.9}{
            \makecell[{{p{\linewidth}}}]{
                \texttt{\tiny{[GM$|$GM]}}
                \texttt{game\_result = WIN} \\
            }
        }
    }
    & & \\ \\

\end{supertabular}
}

\end{document}
